\chapter{Theory and Fundamentals}

The proposed IDS workflow depends on concepts from network security, machine learning, signal representation, spiking neuron models and analog simulation. This chapter explains the technical foundations needed to understand the design choices made in the project. The discussion begins with IDS concepts and dataset structure, then moves to preprocessing, PCA, spike encoding, LIF neuron dynamics and circuit-compatible waveform generation.

\section{Intrusion Detection Systems}

An Intrusion Detection System observes activity in a computing system or network and identifies behaviour that may indicate a security violation. A network-based IDS inspects traffic flows, packet statistics or protocol-level features. It does not necessarily prevent the attack directly; instead, it provides detection, alerting and evidence for response.

IDS methods are commonly classified as:
\begin{enumerate}
\item \textbf{Signature-based detection}: Matches traffic against known malicious signatures. It is reliable for known attacks but weak for unknown variants.
\item \textbf{Anomaly-based detection}: Learns normal behaviour and flags deviations. It can detect new attacks but may produce false positives.
\item \textbf{Machine learning based detection}: Uses labelled or unlabelled data to learn decision boundaries from traffic features.
\end{enumerate}

In a machine learning IDS, each network flow is represented as a feature vector. A classifier then maps the feature vector to a class label such as normal or attack. The quality of detection depends heavily on preprocessing, feature selection, training data and model choice.

\section{UNSW-NB15 Dataset}

The UNSW-NB15 dataset is used as the project dataset. It contains normal traffic and attack traffic generated in a controlled cyber range environment \cite{Moustafa2015}. The dataset includes multiple attack classes such as Fuzzers, Analysis, Backdoors, Denial of Service, Exploits, Generic, Reconnaissance, Shellcode and Worms. It also includes normal network traffic.

The raw dataset contains fields such as protocol, source and destination ports, flow duration, packet counts, byte counts, time-to-live values, TCP-related statistics, service type, state, attack category and label. Some fields are numerical, while others are categorical. This mixture requires careful preprocessing before machine learning or spike encoding can be applied.

\begin{table}[H]
\centering
\caption{Examples of UNSW-NB15 feature groups used in preprocessing}
\label{tab:featuregroups}
\begin{tabular}{|p{0.25\textwidth}|p{0.6\textwidth}|}
\hline
\textbf{Feature group} & \textbf{Examples and significance} \\
\hline
Protocol and service & Protocol, state and service describe the type of network communication. \\
\hline
Traffic volume & Source bytes, destination bytes, source packets and destination packets indicate flow size. \\
\hline
Timing behaviour & Duration and inter-packet timing can indicate scanning, flooding or abnormal communication. \\
\hline
Connection statistics & Count-based features describe repeated connections to hosts, ports or services. \\
\hline
Attack label & Attack category identifies whether the flow is normal or belongs to an attack class. \\
\hline
\end{tabular}
\end{table}

\section{Data Cleaning and Encoding}

Raw network datasets are rarely ready for direct model training. Missing values, placeholder symbols, categorical strings and identifier columns must be handled first. In this project, IP address fields and timestamp fields are removed because they can introduce dataset-specific bias and are not directly suitable for analog neuron mapping. The binary label field is removed when multi-class attack category information is retained.

Categorical features such as protocol and state are converted into numerical form using label encoding. If a categorical feature has $k$ unique values, each value is assigned an integer code. Although label encoding does not preserve semantic distance between categories, it provides a compact representation that can be normalized and processed further. The attack category is also encoded so that the generated samples can preserve class information.

\section{Normalization}

Normalization maps features with different units and ranges into a common numerical interval. Without normalization, high-magnitude features such as byte counts can dominate low-magnitude features such as flags or binary indicators. Min-Max normalization is used:
\begin{equation}
x_{norm} = \frac{x - x_{min}}{x_{max} - x_{min}}
\end{equation}
where $x$ is the original value and $x_{norm}$ is the scaled value. In the hardware context, normalization is also important because voltage-domain spike encoding expects bounded values.

\section{Principal Component Analysis}

Principal Component Analysis reduces the dimensionality of a dataset by transforming the original features into a set of orthogonal components. Each component is a linear combination of original features. The first principal component captures the largest possible variance, the second captures the largest remaining variance orthogonal to the first and so on \cite{Jolliffe2016}.

Let $X$ be the normalized feature matrix. PCA computes directions that maximize projected variance. If $W$ is the matrix containing selected principal component directions, the reduced representation is:
\begin{equation}
Z = XW
\end{equation}
where $Z$ is the reduced feature matrix. In this project, the number of PCA components is selected as eight so that the software feature space maps to eight analog input neurons.

PCA outputs may contain negative values even when the input features are normalized. Therefore, the PCA output is normalized again:
\begin{equation}
z_{scaled} = \frac{z - z_{min}}{z_{max} - z_{min}}
\end{equation}
This produces values suitable for spike encoding.

\section{Spike Encoding}

Spike encoding converts a continuous value into a sequence of binary events. The project implements three encoders.

\subsection{Rate Coding}

In rate coding, a larger input value produces more spikes in a fixed time window. If $x$ is a normalized value and $T$ is the number of timesteps, the number of active spike ticks is:
\begin{equation}
N_{spikes} = \lfloor xT \rfloor
\end{equation}
For example, if $T=10$ and $x=0.7$, approximately seven timesteps are active.

\subsection{Time-to-First-Spike Coding}

In time-to-first-spike coding, the information is represented by the timing of the first spike. Larger values produce earlier spikes:
\begin{equation}
t_{spike} = \lfloor (1 - x)(T - 1) \rfloor
\end{equation}
This coding method is useful when latency is important because strong inputs can be represented quickly.

\subsection{Population Coding}

Population coding represents one feature using multiple neurons or receptive fields. In this project, Gaussian receptive fields are used. For a feature value $x$ and receptive field center $\mu_i$, the activation is:
\begin{equation}
a_i = e^{-\frac{1}{2}\left(\frac{x-\mu_i}{\sigma}\right)^2}
\end{equation}
Each activation $a_i$ is then rate encoded. Population coding expands the representational capacity of the spike input. With eight PCA features and four receptive fields per feature, a sample produces 32 spike trains.

\section{Leaky Integrate-and-Fire Neuron}

The Leaky Integrate-and-Fire neuron models the membrane voltage of a neuron-like element. It integrates incoming input, leaks over time and emits an output spike when the membrane voltage crosses a threshold. A continuous-time LIF model is:
\begin{equation}
\tau_m \frac{dV_m(t)}{dt} = -V_m(t) + RI(t)
\end{equation}
where $V_m(t)$ is membrane voltage, $\tau_m$ is membrane time constant, $R$ is equivalent resistance and $I(t)$ is input current. When:
\begin{equation}
V_m(t) \geq V_{th}
\end{equation}
the neuron fires and the membrane voltage is reset:
\begin{equation}
V_m(t) \leftarrow V_{reset}
\end{equation}

The LIF model is useful in analog design because its blocks correspond naturally to circuit elements. A capacitor can store membrane charge, a resistor can provide leakage, and a comparator can detect threshold crossing. This makes the model suitable for Cadence/LTSpice simulation.

\section{PWL Voltage Sources}

Piecewise Linear voltage sources define arbitrary voltage waveforms using time-voltage pairs. A simulator interpolates between these points. The spike encoder generates PWL text files where 0 V represents no spike and 1 V represents spike activity. A small transition interval is included to avoid ideal vertical edges, which can cause convergence issues in SPICE simulations.

\begin{table}[H]
\centering
\caption{PWL waveform parameters used by the encoder}
\label{tab:pwlparams}
\begin{tabular}{|p{0.35\textwidth}|p{0.35\textwidth}|}
\hline
\textbf{Parameter} & \textbf{Value} \\
\hline
Spike low voltage & 0 V \\
\hline
Spike high voltage & 1 V \\
\hline
Timestep duration & 1 ms \\
\hline
Window length & 10 timesteps \\
\hline
Rise/fall interval & 1\% of timestep \\
\hline
\end{tabular}
\end{table}

\section{Spike Feature Extraction}

After LIF processing, the output waveform can be converted into numerical features. Common features include spike count, firing rate, first spike time, inter-spike interval and output pulse timing. These features provide a bridge between analog neuron output and IDS classification. The same idea can be applied to simulated LIF waveforms or exported Cadence waveforms.

\section{Evaluation Metrics}

Although final result analysis is outside this report scope, the planned IDS evaluation uses standard classification metrics. Accuracy measures total correct predictions. Precision measures how many predicted attacks are truly attacks. Recall measures how many actual attacks are detected. The F1-score combines precision and recall. False positive rate is important in IDS because a high false positive rate can overwhelm security operators.

\section{Software and Simulation Tools}

The software workflow uses Python with NumPy, Pandas and Scikit-learn \cite{Pedregosa2011}. NumPy supports numerical computation, Pandas supports tabular data handling and Scikit-learn supports preprocessing, PCA and classifiers. PyLTSpice is included to support future raw waveform parsing from LTSpice. Cadence or LTSpice is used for circuit-level validation of the LIF neuron.

The fundamentals in this chapter justify the project architecture. The next chapter presents the system design that connects the dataset, spike encoder, PWL interface and LIF circuit.

